% Computational discovery chapter: variation, elliptic integrals, mechanics
\chapter{Variational Problems, Elliptic Integrals, and Mechanics}
\label{ch:discovery_variation_mechanics}

\section{The Brachistochrone}

The brachistochrone asks for the curve along which a particle slides from one point
to another in the least time under gravity. Euler's variational viewpoint replaces
guesswork about curves by a functional whose stationary curve can be studied.

For a cycloid parameterized by

\[
x=a(\theta-\sin\theta),\qquad y=a(1-\cos\theta),
\]

the travel time is $T=\theta\sqrt{a/g}$. The following program finds the parameter
for a prescribed endpoint.

\begin{lstlisting}[language=Python]
from euler import brachistochrone_parameter, brachistochrone_time

theta = brachistochrone_parameter(1.0, 1.0)
print(theta, brachistochrone_time(1.0, 1.0))
\end{lstlisting}

Students can compare this with a straight ramp and a polygonal approximation to see
why the fastest curve is not the shortest curve.

\section{Elliptic Integrals}

The arc length of an ellipse leads to an integral that cannot generally be expressed
using elementary functions. The complete elliptic integrals are

\[
K(k)=\int_0^{\pi/2}\frac{d\theta}{\sqrt{1-k^2\sin^2\theta}},
\qquad
E(k)=\int_0^{\pi/2}\sqrt{1-k^2\sin^2\theta}\,d\theta.
\]

\begin{lstlisting}[language=Python]
from euler import elliptic_k, elliptic_e
from math import pi

print(elliptic_k(0.0), pi/2)
print(elliptic_e(0.0), pi/2)
print(elliptic_k(0.5), elliptic_e(0.5))
\end{lstlisting}

At $k=0$ both integrals reduce to $\pi/2$. This gives an immediate verification of
the numerical method before exploring the genuinely non-elementary cases.

\section{Euler's Equations for a Rigid Body}

For a torque-free body with principal moments $I_1,I_2,I_3$, Euler's equations are

\[
\dot L_1=\frac{I_2-I_3}{I_1}L_2L_3,
\quad
\dot L_2=\frac{I_3-I_1}{I_2}L_3L_1,
\quad
\dot L_3=\frac{I_1-I_2}{I_3}L_1L_2.
\]

\begin{lstlisting}[language=Python]
from euler import integrate_euler_top, rigid_body_energy

moments = (1.0, 2.0, 3.0)
trajectory = integrate_euler_top((1.0, 0.2, 0.1), moments,
                                 step=0.001, steps=5000)
energies = [rigid_body_energy(state, moments) for state in trajectory]
print(max(energies) - min(energies))
\end{lstlisting}

The nearly constant energy is a verification of the physical structure, not merely
of numerical closeness. Changing the step size shows the difference between a
mathematical invariant and numerical drift.

\section{Two-Body Motion}

With units chosen so that the gravitational parameter is $1$, the planar equations
are

\[
\ddot x=-\frac{x}{(x^2+y^2)^{3/2}},\qquad
\ddot y=-\frac{y}{(x^2+y^2)^{3/2}}.
\]

\begin{lstlisting}[language=Python]
from euler import integrate_two_body, two_body_energy

orbit = integrate_two_body((1.0, 0.0, 0.0, 1.0),
                           step=0.001, steps=6283)
print(two_body_energy(orbit[0]), two_body_energy(orbit[-1]))
\end{lstlisting}

The initial circular orbit should return near its starting point after approximately
$2\pi$ time units. Students can vary the initial tangential velocity to produce
elliptic, escape, and collision trajectories.

\section{Exercises}
\begin{exercise}
Compare the cycloid travel time with straight-line motion for the same endpoints.
\end{exercise}
\begin{exercise}
Measure Simpson-rule convergence for $K(0)$ and $E(0)$.
\end{exercise}
\begin{exercise}
Plot the three angular-momentum components of the Euler top and track energy drift.
\end{exercise}
\begin{exercise}
Simulate a non-circular orbit and verify conservation of energy and angular momentum.
\end{exercise}
